Constant weights of magnitude at least \(w_{\min}\) imply the active-weight-region condition on \(B_0\).  The variance-risk absorption lemma at \(d=1\), using \(v_{\min}<v_{\max}\), gives
\[
\mathsf Q_n^{\mu}(s,1;\omega)\asymp n^{-2\alpha_1},
\qquad
\mathsf Q_n^{\sigma}(s,r,1;\omega)\geq c n^{-2\beta_1}.
\tag{15}
\]
The already-established transfer of the Shen--Gao--Witten--Han random-design estimator gives only the upper inequality
\[
\mathsf Q_n^{\sigma}(s,r,1;\omega)\leq Cn^{-2\gamma_1}.
\tag{16}
\]
No lower bound at exponent \(\gamma_1\) is used.

It remains to identify the maximum.  Put
\[
g_1=\frac{4sr}{4sr+2s+r}.
\]
Positive denominators give
\[
g_1<\beta_1
\quad\Longleftrightarrow\quad
4s(2r+1)<4sr+2s+r
\quad\Longleftrightarrow\quad
s<\frac{r}{4r+2}.
\tag{17}
\]
Moreover,
\[
g_1-\alpha_1
=\frac{s\{s(4r-2)+3r\}}
{(4sr+2s+r)(2s+1)}.
\tag{18}
\]
If \(r\geq1/2\), the numerator in (18) is positive.  If \(0<r<1/2\) and \(s<r/(4r+2)\), then
\[
s(2-4r)<\frac{r(2-4r)}{4r+2}<r<3r,
\]
so the numerator is again positive.  Thus \(\gamma_1=\min\{g_1,\beta_1\}\geq\alpha_1\wedge\beta_1\).  Equation (16) is consequently bounded above by a constant times \(n^{-2\alpha_1}\vee n^{-2\beta_1}\).  The two lower bounds in (15) supply the two respective terms.  Their maximum proves the final display.